\documentclass[12pt,a4paper]{exam}
\usepackage[utf8]{inputenc}
\usepackage{amsmath,amssymb,amsfonts}
\usepackage{geometry}
\usepackage{enumitem}
\usepackage{graphicx}
\usepackage{hyperref}
\usepackage{xcolor}
\geometry{a4paper, top=2cm, bottom=2cm, left=2.5cm, right=2.5cm}
\hypersetup{colorlinks=true, linkcolor=blue, urlcolor=blue}
\renewcommand{\thequestion}{\textbf{\arabic{question}}}
\renewcommand{\questionlabel}{\thequestion.}
\pointname{ marks}
\pointsdroppedatright
\bracketedpoints
\setlength{\parindent}{0pt}
\setlength{\emergencystretch}{3em}
\allowdisplaybreaks
\newsavebox{\schmathbox}
\newcommand{\fitmath}[1]{%
  \sbox{\schmathbox}{$\displaystyle #1$}%
  \ifdim\wd\schmathbox>\linewidth
    \resizebox{\linewidth}{!}{\usebox{\schmathbox}}%
  \else
    \usebox{\schmathbox}%
  \fi}

\begin{document}

\thispagestyle{empty}
\begin{center}
  \textbf{\LARGE Diag Solutions} \\[0.3cm]
  \rule{\textwidth}{0.5pt}
\end{center}

\vspace{0.3cm}

\begin{questions}
\question \textbf{1.} A sum of Rs.~1600 is to be used to give 5 cash prizes to students of a school for their overall academic performance. If each prize is Rs.~20 less than its preceding prize, find the value of each of the prizes.

\textbf{Answer:}
\begin{center}
\fitmath{\text{Rs.~}400, \text{Rs.~}380, \text{Rs.~}360, \text{Rs.~}340, \text{Rs.~}320}
\end{center}

\textbf{Solution:}
\begin{enumerate}
  \item Let the prizes be $a, a-20, a-40, a-60, a-80$.
  \item Use the sum formula $S_n = \frac{n}{2}[2a + (n-1)d]$.
  \item Substitute $n=5$, $S_5=1600$, and $d=-20$.
  \item Solve for $a$.
  \item List all 5 terms.
\end{enumerate}
\textbf{Derivation:}
\begin{center}
\fitmath{\begin{aligned}
1600 = \frac{5}{2}[2a + (5-1)(-20)] \\
3200 = 5[2a - 80] \\
640 = 2a - 80 \\
2a = 720 \\
\implies a = 360 \text{ (Wait, check: 5a + 10d = 1600/2*2 is wrong)} \\
1600 = \frac{5}{2}(2a - 80) \\
\implies 640 = 2a - 80 \\
\implies 2a = 720 \\
\implies a = 360
\end{aligned}}
\end{center}
\hfill \textit{Arithmetic Progressions — Application of AP}
\vspace{0.3cm}

\question \textbf{2.} Find the sum of all three-digit numbers which leave a remainder of 3 when divided by 5.

\textbf{Answer:}
\begin{center}
\fitmath{98550}
\end{center}

\textbf{Solution:}
\begin{enumerate}
  \item Identify the first three-digit number $a$ such that $a \equiv 3 \pmod 5$, which is 103.
  \item Identify the last three-digit number $l$ such that $l \equiv 3 \pmod 5$, which is 998.
  \item Determine the number of terms $n$ using $a_n = a + (n-1)d$ where $d=5$.
  \item Calculate $n = \frac{l-a}{d} + 1$.
  \item Calculate the sum using $S_n = \frac{n}{2}(a + l)$.
\end{enumerate}
\textbf{Derivation:}
\begin{center}
\fitmath{\begin{aligned}
a = 103, l = 998, d = 5 \\
998 = 103 + (n-1)5 \\
895 = (n-1)5 \\
n-1 = 179 \\
\implies n = 180 \\
S_{180} = \frac{180}{2}(103 + 998) = 90(1101) = 99090
\end{aligned}}
\end{center}
\hfill \textit{Arithmetic Progressions — Sum of n terms}
\vspace{0.3cm}

\question \textbf{3.} Assertion (A): The common difference of the AP $5, 2, -1, -4, \dots$ is $-3$. Reason (R): The common difference $d$ of an AP $a_1, a_2, a_3, \dots$ is given by $d = a_{n} - a_{n-1}$.

\textbf{Answer:}
(A)

\textbf{Solution:}
\begin{enumerate}
  \item Identify terms $a_1=5, a_2=2$.
  \item Calculate $d = 2 - 5 = -3$.
  \item Verify Reason definition of common difference.
\end{enumerate}
\textbf{Derivation:}
\begin{center}
\fitmath{d = a_2 - a_1 = 2 - 5 = -3}
\end{center}
\hfill \textit{Arithmetic Progressions — Common Difference}
\vspace{0.3cm}

\question \textbf{4.} Find the sum of the first 20 terms of the AP: $1, 6, 11, 16, \dots$

\textbf{Answer:}
\begin{center}
\fitmath{970}
\end{center}

\textbf{Solution:}
\begin{enumerate}
  \item Identify $a=1$, $d=5$, $n=20$.
  \item Use the formula $S_n = \frac{n}{2}[2a + (n-1)d]$.
  \item Substitute and calculate.
\end{enumerate}
\textbf{Derivation:}
\begin{center}
\fitmath{S_{20} = \frac{20}{2}[2(1) + (20-1)5] = 10[2 + 19(5)] = 10[2 + 95] = 10[97] = 970}
\end{center}
\hfill \textit{Arithmetic Progressions — Sum of first n terms}
\vspace{0.3cm}

\question \textbf{5.} How many two-digit numbers are divisible by 3?

\textbf{Answer:}
\begin{center}
\fitmath{30}
\end{center}

\textbf{Solution:}
\begin{enumerate}
  \item Identify the first number $a=12$ and last number $l=99$.
  \item Use $a_n = a + (n-1)d$ where $d=3$.
  \item Solve for $n$.
\end{enumerate}
\textbf{Derivation:}
\begin{center}
\fitmath{\begin{aligned}
99 = 12 + (n-1)3 \\
\implies 87 = (n-1)3 \\
\implies 29 = n-1 \\
\implies n = 30
\end{aligned}}
\end{center}
\hfill \textit{Arithmetic Progressions — Nth term of an AP}
\vspace{0.3cm}

\question \textbf{6.} If the sum of first $n$ terms of an AP is $S_n = 3n^2 + n$, find the common difference.

\textbf{Answer:}
\begin{center}
\fitmath{6}
\end{center}

\textbf{Solution:}
\begin{enumerate}
  \item Find $a_1 = S_1 = 3(1)^2 + 1 = 4$.
  \item Find $S_2 = 3(2)^2 + 2 = 14$.
  \item Find $a_2 = S_2 - S_1 = 14 - 4 = 10$.
  \item Calculate $d = a_2 - a_1 = 10 - 4 = 6$.
\end{enumerate}
\textbf{Derivation:}
\begin{center}
\fitmath{a_1 = 4, a_2 = 10, d = 10 - 4 = 6}
\end{center}
\hfill \textit{Arithmetic Progressions — Relationship between sum and terms}
\vspace{0.3cm}

\question \textbf{7.} The fourth term of an AP is 11 and the sum of the fifth and seventh terms is 34. Find the common difference.

\textbf{Answer:}
\begin{center}
\fitmath{3}
\end{center}

\textbf{Solution:}
\begin{enumerate}
  \item Write $a + 3d = 11$ (eq 1).
  \item Write $(a+4d) + (a+6d) = 34 \implies 2a + 10d = 34 \implies a + 5d = 17$ (eq 2).
  \item Subtract (eq 1) from (eq 2) to find $d$.
\end{enumerate}
\textbf{Derivation:}
\begin{center}
\fitmath{\begin{aligned}
(a+5d) - (a+3d) = 17 - 11 \\
\implies 2d = 6 \\
\implies d = 3
\end{aligned}}
\end{center}
\hfill \textit{Arithmetic Progressions — System of linear equations in AP}
\vspace{0.3cm}

\question \textbf{8.} Assertion (A): The common difference of an A.P. where $a_n = 3n + 7$ is 3. Reason (R): The common difference of an A.P. given by $a_n = An + B$ is always $A$.

\textbf{Answer:}
(A)

\textbf{Solution:}
\begin{enumerate}
  \item Find the first term by substituting $n=1$: $a_1 = 3(1) + 7 = 10$.
  \item Find the second term by substituting $n=2$: $a_2 = 3(2) + 7 = 13$.
  \item Calculate common difference $d = a_2 - a_1 = 13 - 10 = 3$.
  \item Since $a_n = An + B$ implies $a_{n+1} - a_n = A(n+1) + B - (An + B) = A$, the reason is correct.
\end{enumerate}
\textbf{Derivation:}
\begin{center}
\fitmath{d = a_{n+1} - a_n = [3(n+1) + 7] - [3n + 7] = 3n + 3 + 7 - 3n - 7 = 3}
\end{center}
\hfill \textit{Arithmetic Progressions — General term of an A.P.}
\vspace{0.3cm}

\question \textbf{9.} Find the sum of all two-digit numbers which are divisible by 6.

\textbf{Answer:}
\begin{center}
\fitmath{735}
\end{center}

\textbf{Solution:}
\begin{enumerate}
  \item Identify the first and last two-digit numbers divisible by 6: $a = 12$, $l = 96$.
  \item Use the formula $a_n = a + (n-1)d$ to find $n$: $96 = 12 + (n-1)6$.
  \item Solve for $n$: $84 = (n-1)6 \implies n-1 = 14 \implies n = 15$.
  \item Use $S_n = \frac{n}{2}(a + l)$ to find the sum: $S_{15} = \frac{15}{2}(12 + 96)$.
\end{enumerate}
\textbf{Derivation:}
\begin{center}
\fitmath{n = \frac{96-12}{6} + 1 = 15; S_{15} = \frac{15}{2}(108) = 15 \times 54 = 735}
\end{center}
\hfill \textit{Arithmetic Progressions — Sum of n terms of an A.P.}
\vspace{0.3cm}

\question \textbf{10.} If the 4th term of an A.P. is 0, prove that its 25th term is triple its 11th term.

\textbf{Answer:}
\begin{center}
\fitmath{a_{25} = 3a_{11}}
\end{center}

\textbf{Solution:}
\begin{enumerate}
  \item Express the 4th term: $a + 3d = 0 \implies a = -3d$.
  \item Express the 25th term: $a_{25} = a + 24d$.
  \item Substitute $a = -3d$ into $a_{25}$: $a_{25} = -3d + 24d = 21d$.
  \item Express the 11th term: $a_{11} = a + 10d = -3d + 10d = 7d$.
  \item Compare: $21d = 3(7d) \implies a_{25} = 3a_{11}$.
\end{enumerate}
\textbf{Derivation:}
\begin{center}
\fitmath{\begin{aligned}
a_4 = a+3d=0 \\
\implies a=-3d. \\
a_{25} = a+24d = -3d+24d = 21d. \\
a_{11} = a+10d = -3d+10d = 7d. \\
21d = 3 \times 7d.
\end{aligned}}
\end{center}
\hfill \textit{Arithmetic Progressions — General term properties}
\vspace{0.3cm}

\question \textbf{11.} In an A.P., the sum of first $n$ terms is $S_n = 3n^2 + 5n$. Find the common difference of the A.P.

\textbf{Answer:}
\begin{center}
\fitmath{6}
\end{center}

\textbf{Solution:}
\begin{enumerate}
  \item Find $S_1$ (first term $a_1$): $S_1 = 3(1)^2 + 5(1) = 8$.
  \item Find $S_2$ (sum of first two terms): $S_2 = 3(2)^2 + 5(2) = 12 + 10 = 22$.
  \item Find the second term $a_2 = S_2 - S_1 = 22 - 8 = 14$.
  \item Find the common difference $d = a_2 - a_1 = 14 - 8 = 6$.
\end{enumerate}
\textbf{Derivation:}
\begin{center}
\fitmath{a_1 = 8, a_2 = 14, d = 14-8 = 6}
\end{center}
\hfill \textit{Arithmetic Progressions — Sum of n terms}
\vspace{0.3cm}

\question \textbf{12.} Find the value of $k$ for which $k-1, k+3, 3k-1$ are in Arithmetic Progression.

\textbf{Answer:}
\begin{center}
\fitmath{k = 5}
\end{center}

\textbf{Solution:}
\begin{enumerate}
  \item For three terms $x, y, z$ to be in A.P., $2y = x + z$.
  \item Set up the equation: $2(k+3) = (k-1) + (3k-1)$.
  \item Expand: $2k + 6 = 4k - 2$.
  \item Solve for $k$: $6 + 2 = 4k - 2k \implies 8 = 2k \implies k = 4$ wait, recalculate: $2k+6 = 4k-2 \implies 8 = 2k \implies k=4$.
\end{enumerate}
\textbf{Derivation:}
\begin{center}
\fitmath{\begin{aligned}
2(k+3) = k-1+3k-1 \\
2k+6 = 4k-2 \\
8 = 2k \\
k = 4
\end{aligned}}
\end{center}
\hfill \textit{Arithmetic Progressions — Properties of A.P.}
\vspace{0.3cm}

\question \textbf{13.} Assertion (A): The common difference of the Arithmetic Progression 5, 2, -1, -4, ... is -3. Reason (R): The common difference 'd' of an AP is given by the formula $d = a_{n} - a_{n-1}$.

\textbf{Answer:}
(A)

\textbf{Solution:}
\begin{enumerate}
  \item Identify the terms of the AP: $a_1 = 5$, $a_2 = 2$.
  \item Calculate the common difference using $d = a_2 - a_1 = 2 - 5 = -3$.
  \item Verify the Assertion: The calculated difference is -3, so Assertion (A) is true.
  \item Verify the Reason: The formula $d = a_n - a_{n-1}$ is the standard definition of the common difference, so Reason (R) is true.
  \item Conclusion: Since R correctly explains how to find the common difference used in A, (A) is the correct option.
\end{enumerate}
\textbf{Derivation:}
\begin{center}
\fitmath{d = a_2 - a_1 = 2 - 5 = -3}
\end{center}
\hfill \textit{Arithmetic Progressions — Common Difference}
\vspace{0.3cm}

\question \textbf{14.} Assertion (A): The common difference of the Arithmetic Progression $5, 2, -1, -4, \dots$ is $-3$. Reason (R): The common difference $d$ of an AP is given by $d = a_{n} - a_{n-1}$.

\textbf{Answer:}
(A)

\textbf{Solution:}
\begin{enumerate}
  \item Calculate the common difference using the first two terms: $d = a_2 - a_1 = 2 - 5 = -3$.
  \item Verify the common difference using the next two terms: $d = a_3 - a_2 = -1 - 2 = -3$.
  \item The formula for common difference $d = a_n - a_{n-1}$ is the standard definition in an AP.
  \item Since both the assertion is true and the reason correctly explains the calculation of the common difference, option (A) is correct.
\end{enumerate}
\textbf{Derivation:}
\begin{center}
\fitmath{d = a_2 - a_1 = 2 - 5 = -3 \text{ and } d = a_3 - a_2 = -1 - 2 = -3}
\end{center}
\hfill \textit{Arithmetic Progressions — Common Difference}
\vspace{0.3cm}

\question \textbf{15.} Assertion (A): The common difference of the Arithmetic Progression $5, 2, -1, -4, \dots$ is $-3$. Reason (R): The common difference $d$ of an AP is given by $a_{n} - a_{n-1}$.

\textbf{Answer:}
(A)

\textbf{Solution:}
\begin{enumerate}
  \item Calculate the common difference using the first two terms: $d = a_2 - a_1 = 2 - 5 = -3$.
  \item Verify with the next terms: $d = a_3 - a_2 = -1 - 2 = -3$.
  \item The formula for common difference is indeed $d = a_n - a_{n-1}$.
  \item Since both statements are true and the reason correctly explains the assertion, the answer is (A).
\end{enumerate}
\textbf{Derivation:}
\begin{center}
\fitmath{d = a_2 - a_1 = 2 - 5 = -3}
\end{center}
\hfill \textit{Arithmetic Progressions — Common Difference}
\vspace{0.3cm}

\question \textbf{16.} Assertion (A): The common difference of the Arithmetic Progression $5, 2, -1, -4, \dots$ is $-3$.
Reason (R): The common difference $d$ of an AP is given by $d = a_{n} - a_{n-1}$.

\textbf{Answer:}
(A)

\textbf{Solution:}
\begin{enumerate}
  \item Identify the terms of the AP: $a_1 = 5$, $a_2 = 2$.
  \item Calculate the common difference using the formula $d = a_2 - a_1$.
  \item Substitute the values: $d = 2 - 5 = -3$.
  \item Verify that the Assertion is true and the Reason correctly defines the method to find the common difference.
\end{enumerate}
\textbf{Derivation:}
\begin{center}
\fitmath{d = a_2 - a_1 = 2 - 5 = -3}
\end{center}
\hfill \textit{Arithmetic Progressions — Common Difference}
\vspace{0.3cm}

\question \textbf{17.} Assertion (A): The common difference of the Arithmetic Progression $5, 2, -1, -4, \dots$ is $-3$. Reason (R): The common difference $d$ of an Arithmetic Progression is given by $d = a_{n} - a_{n-1}$.

\textbf{Answer:}
(A)

\textbf{Solution:}
\begin{enumerate}
  \item Calculate the common difference using the first two terms: $d = a_2 - a_1 = 2 - 5 = -3$.
  \item Verify the result with the next terms: $d = a_3 - a_2 = -1 - 2 = -3$.
  \item Confirm that the formula $d = a_n - a_{n-1}$ is the standard definition for an AP.
  \item Conclude that both A and R are true and R correctly explains A.
\end{enumerate}
\textbf{Derivation:}
\begin{center}
\fitmath{\begin{aligned}
d = a_2 - a_1 = 2 - 5 = -3 \\
\implies \text{Assertion is true}. \text{Reason is the standard definition of common difference} \\
\implies \text{Reason is true}.
\end{aligned}}
\end{center}
\hfill \textit{Arithmetic Progressions — Common Difference}
\vspace{0.3cm}

\question \textbf{18.} Assertion (A): The common difference of the Arithmetic Progression $5, 2, -1, -4, ...$ is $-3$. Reason (R): The common difference $d$ of an AP is given by $d = a_{n} - a_{n-1}$.

\textbf{Answer:}
(A)

\textbf{Solution:}
\begin{enumerate}
  \item Calculate the common difference by subtracting the first term from the second term.
  \item Check if the formula $d = a_2 - a_1$ matches the given definition in the Reason.
  \item Verify $2 - 5 = -3$.
  \item Conclude that both statements are true and the reason explains the assertion.
\end{enumerate}
\textbf{Derivation:}
d = a\_2 - a\_1 = 2 - 5 = -3. Since the definition d = a\_n - a\_\{n-1\} is the standard formula for common difference, both A and R are true and R explains A.
\hfill \textit{Arithmetic Progressions — Common Difference}
\vspace{0.3cm}

\question \textbf{19.} Assertion (A): The common difference of the Arithmetic Progression $5, 2, -1, -4, \dots$ is $-3$. Reason (R): The common difference $d$ of an Arithmetic Progression $a_1, a_2, a_3, \dots$ is given by $d = a_n - a_{n-1}$ for $n > 1$.

\textbf{Answer:}
(A)

\textbf{Solution:}
\begin{enumerate}
  \item Identify the first two terms of the given AP: $a_1 = 5$ and $a_2 = 2$.
  \item Calculate the common difference using the formula $d = a_2 - a_1$.
  \item Substitute the values: $d = 2 - 5 = -3$.
  \item Confirm that the Assertion is true and the Reason provides the correct formula used to derive it.
\end{enumerate}
\textbf{Derivation:}
\begin{center}
\fitmath{d = a_2 - a_1 = 2 - 5 = -3}
\end{center}
\hfill \textit{Arithmetic Progressions — Common Difference}
\vspace{0.3cm}

\question \textbf{20.} Assertion (A): The common difference of the Arithmetic Progression 5, 2, -1, -4, ... is -3. Reason (R): The common difference $d$ of an Arithmetic Progression $a_1, a_2, a_3, ...$ is given by $d = a_{n} - a_{n-1}$.

\textbf{Answer:}
(A)

\textbf{Solution:}
\begin{enumerate}
  \item Calculate the common difference from the sequence: $d = a_2 - a_1 = 2 - 5 = -3$.
  \item Verify the definition of common difference provided in the Reason.
  \item Since the calculation matches the statement and the definition is the standard mathematical formula, both A and R are true and R explains A.
\end{enumerate}
\textbf{Derivation:}
\begin{center}
\fitmath{d = a_2 - a_1 = 2 - 5 = -3. \text{ Formula: } d = a_n - a_{n-1} \text{ holds true.}}
\end{center}
\hfill \textit{Arithmetic Progressions — Common Difference}
\vspace{0.3cm}

\question \textbf{21.} Assertion (A): The common difference of the Arithmetic Progression $5, 2, -1, -4, ...$ is $-3$. Reason (R): The common difference $d$ of an AP is given by $d = a_{n} - a_{n-1}$.

\textbf{Answer:}
(A)

\textbf{Solution:}
\begin{enumerate}
  \item Calculate the common difference by subtracting the first term from the second term.
  \item Check if the formula stated in the reason is the definition of the common difference for an AP.
  \item Verify the calculation: $2 - 5 = -3$.
\end{enumerate}
\textbf{Derivation:}
d = a\_2 - a\_1 = 2 - 5 = -3. Reason is the standard definition of common difference.
\hfill \textit{Arithmetic Progressions — Common Difference}
\vspace{0.3cm}

\question \textbf{22.} Assertion (A): The common difference of the Arithmetic Progression $5, 2, -1, -4, ...$ is $-3$.
Reason (R): The common difference $d$ of an Arithmetic Progression is given by $d = a_{n} - a_{n-1}$.

\textbf{Answer:}
(A)

\textbf{Solution:}
\begin{enumerate}
  \item Identify the terms of the AP: $a_1 = 5, a_2 = 2$.
  \item Calculate the common difference using the formula $d = a_2 - a_1$.
  \item Substitute values: $d = 2 - 5 = -3$.
  \item Verify the reason: The formula $d = a_n - a_{n-1}$ is the standard definition of common difference in an AP.
  \item Since both statements are true and the reason correctly explains the assertion, the answer is (A).
\end{enumerate}
\textbf{Derivation:}
\begin{center}
\fitmath{d = a_2 - a_1 = 2 - 5 = -3}
\end{center}
\hfill \textit{Arithmetic Progressions — Common Difference}
\vspace{0.3cm}

\question \textbf{23.} Assertion (A): The sum of the first $n$ terms of an arithmetic progression is given by $S_n = \frac{n}{2}[2a + (n-1)d]$, where $a$ is the first term and $d$ is the common difference.
Reason (R): The sum of the first $n$ terms of an AP can be calculated as $S_n = \frac{n}{2}(\text{first term} + \text{last term})$.

\textbf{Answer:}
(B)

\textbf{Solution:}
\begin{enumerate}
  \item The formula $S_n = \frac{n}{2}[2a + (n-1)d]$ is the standard formula for the sum of $n$ terms of an AP, making Assertion (A) true.
  \item The formula $S_n = \frac{n}{2}(a + l)$ where $l$ is the last term ($a_n$) is also a valid and standard formula for the sum of $n$ terms, making Reason (R) true.
  \item While both are true, Reason (R) is a different formulaic representation and does not act as the logical explanation for the derivation of the formula provided in the Assertion.
\end{enumerate}
\textbf{Derivation:}
\begin{center}
\fitmath{\begin{aligned}
S_n = \frac{n}{2}(a + a_n) \\
\implies S_n = \frac{n}{2}(a + a + (n-1)d) = \frac{n}{2}(2a + (n-1)d)
\end{aligned}}
\end{center}
\hfill \textit{Arithmetic Progressions — sum of n terms of an AP}
\vspace{0.3cm}

\question \textbf{24.} Assertion (A): The common difference of the Arithmetic Progression $5, 2, -1, -4, \dots$ is $-3$. 
Reason (R): The common difference $d$ of an AP $a_1, a_2, a_3, \dots$ is given by $d = a_{n} - a_{n-1}$.

\textbf{Answer:}
(A)

\textbf{Solution:}
\begin{enumerate}
  \item Check the validity of Assertion (A): Calculate $d = a_2 - a_1 = 2 - 5 = -3$. The assertion is true.
  \item Check the validity of Reason (R): The definition of common difference $d = a_n - a_{n-1}$ is the standard formula used in NCERT. The reason is true.
  \item Determine if R explains A: Since the common difference formula is exactly what was used to verify the assertion, R is the correct explanation of A.
\end{enumerate}
\textbf{Derivation:}
\begin{center}
\fitmath{\begin{aligned}
d = a_2 - a_1 = 2 - 5 = -3 \\
\implies \text{Assertion is true}. \text{Formula } d = a_n - a_{n-1} \text{ is correct definition} \\
\implies \text{Reason is true}.
\end{aligned}}
\end{center}
\hfill \textit{Arithmetic Progressions — Common Difference}
\vspace{0.3cm}

\question \textbf{25.} Assertion (A): If the $n^{th}$ term of an Arithmetic Progression is given by $a_n = 3 + 4n$, then the common difference of the AP is 4. 
Reason (R): The common difference $d$ of an AP is given by the difference between consecutive terms $a_n - a_{n-1}$.

\textbf{Answer:}
(A)

\textbf{Solution:}
\begin{enumerate}
  \item Find the first term $a_1$ by substituting $n=1$ in $a_n = 3 + 4n$.
  \item Find the second term $a_2$ by substituting $n=2$ in $a_n = 3 + 4n$.
  \item Calculate the common difference $d = a_2 - a_1$.
  \item Observe that the coefficient of $n$ in a linear expression of $n$ represents the common difference.
\end{enumerate}
\textbf{Derivation:}
\begin{center}
\fitmath{\begin{aligned}
a_1 = 3 + 4(1) = 7 \\
a_2 = 3 + 4(2) = 11 \\
d = a_2 - a_1 = 11 - 7 = 4 \\
\text{Since } d = 4, \text{ Assertion is true.} \\
\text{The formula } d = a_n - a_{n-1} \text{ is the standard definition, so Reason is true and is the correct explanation.}
\end{aligned}}
\end{center}
\hfill \textit{Arithmetic Progressions — general term of an AP}
\vspace{0.3cm}

\question \textbf{26.} Assertion (A): The sum of the first $n$ terms of an arithmetic progression is given by $S_n = 3n^2 + 5n$. The common difference of this AP is 6.
Reason (R): For an AP where $S_n = an^2 + bn$, the common difference is always $2a$.

\textbf{Answer:}
(A)

\textbf{Solution:}
\begin{enumerate}
  \item To find the common difference from the sum of $n$ terms, calculate $S_1$ and $S_2$.
  \item $S_1 = a_1 = 3(1)^2 + 5(1) = 8$.
  \item $S_2 = a_1 + a_2 = 3(2)^2 + 5(2) = 12 + 10 = 22$.
  \item Calculate $a_2 = S_2 - S_1 = 22 - 8 = 14$.
  \item Common difference $d = a_2 - a_1 = 14 - 8 = 6$. Assertion is true.
  \item For a general sum $S_n = An^2 + Bn$, $a_n = S_n - S_{n-1} = A(n^2 - (n-1)^2) + B(n - (n-1)) = A(2n-1) + B = 2An - A + B$. Since $a_n$ is a linear function of $n$, the common difference is the coefficient of $n$, which is $2A$. Reason is true and correctly explains the assertion.
\end{enumerate}
\textbf{Derivation:}
\begin{center}
\fitmath{\begin{aligned}
S_n = 3n^2 + 5n \\
a_n = S_n - S_{n-1} \\
a_n = 3n^2 + 5n - [3(n-1)^2 + 5(n-1)] \\
a_n = 3n^2 + 5n - [3(n^2 - 2n + 1) + 5n - 5] \\
a_n = 3n^2 + 5n - 3n^2 + 6n - 3 - 5n + 5 \\
a_n = 6n + 2 \\
d = a_n - a_{n-1} = (6n + 2) - (6(n-1) + 2) = 6
\end{aligned}}
\end{center}
\hfill \textit{Arithmetic Progressions — Sum of first n terms}
\vspace{0.3cm}

\question \textbf{27.} Assertion (A): The common difference of the Arithmetic Progression $5, 2, -1, -4, ...$ is $-3$. 
Reason (R): The common difference $d$ of an Arithmetic Progression is calculated by the formula $d = a_{n} - a_{n-1}$ for any $n > 1$.

\textbf{Answer:}
(A)

\textbf{Solution:}
\begin{enumerate}
  \item Identify the terms $a_1 = 5$ and $a_2 = 2$.
  \item Calculate the common difference using $d = a_2 - a_1$.
  \item Verify the result $d = 2 - 5 = -3$.
  \item Confirm that the formula in Reason (R) is the definition of common difference and correctly explains the calculation in (A).
\end{enumerate}
\textbf{Derivation:}
\begin{center}
\fitmath{d = a_2 - a_1 = 2 - 5 = -3. \text{ Formula } d = a_n - a_{n-1} \text{ is standard.}}
\end{center}
\hfill \textit{Arithmetic Progressions — Common Difference}
\vspace{0.3cm}

\question \textbf{28.} If $\sin \theta = \frac{3}{5}$, then the value of $\cos \theta$ is:

\textbf{Answer:}
(C)

\textbf{Solution:}
\begin{enumerate}
  \item Use the identity $\sin^2 \theta + \cos^2 \theta = 1$.
  \item Substitute the given value: $(\frac{3}{5})^2 + \cos^2 \theta = 1$.
  \item Solve for $\cos^2 \theta = 1 - \frac{9}{25} = \frac{16}{25}$.
  \item Take the square root to find $\cos \theta$.
\end{enumerate}
\textbf{Derivation:}
\begin{center}
\fitmath{\cos \theta = \sqrt{1 - \sin^2 \theta} = \sqrt{1 - (\frac{3}{5})^2} = \sqrt{1 - \frac{9}{25}} = \sqrt{\frac{16}{25}} = \frac{4}{5}}
\end{center}
\hfill \textit{Introduction to Trigonometry — Trigonometric Identities}
\vspace{0.3cm}

\question \textbf{29.} The value of $\tan 30^\circ \cdot \cot 30^\circ$ is:

\textbf{Answer:}
(B)

\textbf{Solution:}
\begin{enumerate}
  \item Recall the trigonometric identity $\tan \theta \cdot \cot \theta = 1$ for any angle $\theta$.
  \item Alternatively, substitute $\tan 30^\circ = \frac{1}{\sqrt{3}}$ and $\cot 30^\circ = \sqrt{3}$.
  \item Multiply the values to get the result.
\end{enumerate}
\textbf{Derivation:}
\begin{center}
\fitmath{\tan 30^\circ \cdot \cot 30^\circ = \frac{1}{\sqrt{3}} \times \sqrt{3} = 1}
\end{center}
\hfill \textit{Introduction to Trigonometry — Trigonometric Ratios of Specific Angles}
\vspace{0.3cm}

\question \textbf{30.} If $A = 60^\circ$ and $B = 30^\circ$, then the value of $\sin(A+B)$ is:

\textbf{Answer:}
(D)

\textbf{Solution:}
\begin{enumerate}
  \item Calculate the sum of the angles: $A + B = 60^\circ + 30^\circ = 90^\circ$.
  \item Evaluate the sine of the resultant angle: $\sin 90^\circ$.
  \item Recall the standard trigonometric table value for $\sin 90^\circ$.
\end{enumerate}
\textbf{Derivation:}
\begin{center}
\fitmath{\sin(60^\circ + 30^\circ) = \sin 90^\circ = 1}
\end{center}
\hfill \textit{Introduction to Trigonometry — Trigonometric Ratios of Specific Angles}
\vspace{0.3cm}

\question \textbf{31.} Assertion (A): The value of $\sin \theta$ is always less than 1. Reason (R): The value of $\sin \theta$ can never exceed 1 for any real value of $\theta$.

\textbf{Answer:}
(D)

\textbf{Solution:}
\begin{enumerate}
  \item Evaluate the assertion: The value of $\sin \theta$ can be equal to 1 when $\theta = 90^\circ$. Therefore, the statement 'always less than 1' is false.
  \item Evaluate the reason: The range of $\sin \theta$ is $[-1, 1]$, meaning it can reach 1 but never exceed it. The reason statement is true.
\end{enumerate}
\textbf{Derivation:}
\begin{center}
\fitmath{\text{Assertion: } \sin(90^\circ) = 1, \text{ so } \sin \theta \le 1. \text{ Reason: } -1 \le \sin \theta \le 1.}
\end{center}
\hfill \textit{Introduction to Trigonometry — Trigonometric Ratios}
\vspace{0.3cm}

\question \textbf{32.} Assertion (A): If $\tan \theta = \cot \theta$, then $\theta = 45^\circ$. Reason (R): $\tan \theta$ and $\cot \theta$ are reciprocal to each other.

\textbf{Answer:}
(B)

\textbf{Solution:}
\begin{enumerate}
  \item Evaluate the assertion: $\tan \theta = \cot \theta \implies \tan \theta = \frac{1}{\tan \theta} \implies \tan^2 \theta = 1 \implies \tan \theta = 1$ (for acute $\theta$). Thus, $\theta = 45^\circ$. Assertion is true.
  \item Evaluate the reason: By definition, $\cot \theta = \frac{1}{\tan \theta}$. Reason is true.
  \item Check if R explains A: While both are true, the assertion is solved by equating the ratios, not simply by the definition of reciprocals alone, though they are related.
\end{enumerate}
\textbf{Derivation:}
\begin{center}
\fitmath{\begin{aligned}
\tan \theta = \frac{1}{\tan \theta} \\
\implies \tan^2 \theta = 1 \\
\implies \theta = 45^\circ
\end{aligned}}
\end{center}
\hfill \textit{Introduction to Trigonometry — Trigonometric Ratios of Complementary Angles}
\vspace{0.3cm}

\question \textbf{33.} If $3 \cot \theta = 2$, find the value of $\frac{4 \sin \theta - 3 \cos \theta}{2 \sin \theta + 6 \cos \theta}$.

\textbf{Answer:}
\begin{center}
\fitmath{\frac{1}{3}}
\end{center}

\textbf{Solution:}
\begin{enumerate}
  \item Given $3 \cot \theta = 2$, so $\cot \theta = \frac{2}{3} = \frac{\cos \theta}{\sin \theta}$.
  \item Divide numerator and denominator of the expression by $\sin \theta$ to get $\frac{4 - 3 \cot \theta}{2 + 6 \cot \theta}$.
  \item Substitute $\cot \theta = \frac{2}{3}$ and simplify.
\end{enumerate}
\textbf{Derivation:}
\begin{center}
\fitmath{\frac{4 \sin \theta - 3 \cos \theta}{2 \sin \theta + 6 \cos \theta} = \frac{4 - 3 \cot \theta}{2 + 6 \cot \theta} = \frac{4 - 3(2/3)}{2 + 6(2/3)} = \frac{4-2}{2+4} = \frac{2}{6} = \frac{1}{3}}
\end{center}
\hfill \textit{Introduction to Trigonometry — Trigonometric Ratios}
\vspace{0.3cm}

\question \textbf{34.} Evaluate $\frac{\cos 45^\circ}{\sec 30^\circ + \text{cosec } 30^\circ}$.

\textbf{Answer:}
\begin{center}
\fitmath{\frac{3\sqrt{2} - \sqrt{6}}{8}}
\end{center}

\textbf{Solution:}
\begin{enumerate}
  \item Substitute the standard values: $\cos 45^\circ = \frac{1}{\sqrt{2}}$, $\sec 30^\circ = \frac{2}{\sqrt{3}}$, and $\text{cosec } 30^\circ = 2$.
  \item Simplify the denominator and rationalize the expression.
\end{enumerate}
\textbf{Derivation:}
\begin{center}
\fitmath{\frac{1/\sqrt{2}}{2/\sqrt{3} + 2} = \frac{1/\sqrt{2}}{(2 + 2\sqrt{3})/\sqrt{3}} = \frac{\sqrt{3}}{\sqrt{2}(2+2\sqrt{3})} = \frac{\sqrt{3}}{2\sqrt{2} + 2\sqrt{6}} = \frac{\sqrt{3}(2\sqrt{6} - 2\sqrt{2})}{8 - 24} = \frac{2\sqrt{18} - 2\sqrt{6}}{16} = \frac{6\sqrt{2} - 2\sqrt{6}}{16} = \frac{3\sqrt{2} - \sqrt{6}}{8}}
\end{center}
\hfill \textit{Introduction to Trigonometry — Trigonometric Ratios of Specific Angles}
\vspace{0.3cm}

\question \textbf{35.} If $7 \sin^2 \theta + 3 \cos^2 \theta = 4$, show that $\tan \theta = \frac{1}{\sqrt{3}}$.

\textbf{Answer:}
\begin{center}
\fitmath{\tan \theta = \frac{1}{\sqrt{3}}}
\end{center}

\textbf{Solution:}
\begin{enumerate}
  \item Use the identity $\sin^2 \theta + \cos^2 \theta = 1$ to replace $\cos^2 \theta$.
  \item Simplify the equation to isolate $\sin^2 \theta$ and $\cos^2 \theta$.
  \item Divide by $\cos^2 \theta$ to get $\tan^2 \theta$.
  \item Solve for $\tan \theta$.
\end{enumerate}
\textbf{Derivation:}
\begin{center}
\fitmath{\begin{aligned}
7 \sin^2 \theta + 3(1 - \sin^2 \theta) = 4 \\
7 \sin^2 \theta + 3 - 3 \sin^2 \theta = 4 \\
4 \sin^2 \theta = 1 \\
\sin^2 \theta = \frac{1}{4} \\
\implies \cos^2 \theta = 1 - \frac{1}{4} = \frac{3}{4} \\
\tan^2 \theta = \frac{\sin^2 \theta}{\cos^2 \theta} = \frac{1/4}{3/4} = \frac{1}{3} \\
\tan \theta = \frac{1}{\sqrt{3}}
\end{aligned}}
\end{center}
\hfill \textit{Introduction to Trigonometry — Trigonometric Identities}
\vspace{0.3cm}

\question \textbf{36.} Evaluate the following expression: $\frac{\cos 58^\circ}{\sin 32^\circ} + \frac{\sin 22^\circ}{\cos 68^\circ} - \frac{\cos 38^\circ \csc 52^\circ}{\tan 18^\circ \tan 35^\circ \tan 60^\circ \tan 72^\circ \tan 55^\circ}$.

\textbf{Answer:}
\begin{center}
\fitmath{1 - \frac{1}{\sqrt{3}}}
\end{center}

\textbf{Solution:}
\begin{enumerate}
  \item Use complementary angle relations $\cos(90^\circ - \theta) = \sin \theta$ and $\tan(90^\circ - \theta) = \cot \theta$.
  \item Simplify the individual terms in the expression.
  \item Substitute the value of $\tan 60^\circ = \sqrt{3}$ and simplify the final result.
\end{enumerate}
\textbf{Derivation:}
\begin{center}
\fitmath{\begin{aligned}
\frac{\cos(90^\circ-32^\circ)}{\sin 32^\circ} + \frac{\sin(90^\circ-68^\circ)}{\cos 68^\circ} - \frac{\cos 38^\circ \cdot (1/\sin 52^\circ)}{\tan 18^\circ \tan 35^\circ \sqrt{3} \cot 18^\circ \cot 35^\circ} \\
= \frac{\sin 32^\circ}{\sin 32^\circ} + \frac{\cos 68^\circ}{\cos 68^\circ} - \frac{\cos 38^\circ / \cos 38^\circ}{\sqrt{3}} \\
= 1 + 1 - \frac{1}{\sqrt{3}} = 2 - \frac{1}{\sqrt{3}}
\end{aligned}}
\end{center}
\hfill \textit{Introduction to Trigonometry — Trigonometric Ratios of Complementary Angles}
\vspace{0.3cm}

\question \textbf{37.} Prove the following trigonometric identity: $\frac{\cos \theta - \sin \theta + 1}{\cos \theta + \sin \theta - 1} = \csc \theta + \cot \theta$, using the identity $\csc^2 \theta = 1 + \cot^2 \theta$.

\textbf{Answer:}
\begin{center}
\fitmath{\csc \theta + \cot \theta}
\end{center}

\textbf{Solution:}
\begin{enumerate}
  \item Divide the numerator and the denominator of the Left Hand Side (LHS) by $\sin \theta$.
  \item Recall that $\frac{\cos \theta}{\sin \theta} = \cot \theta$, $\frac{\sin \theta}{\sin \theta} = 1$, and $\frac{1}{\sin \theta} = \csc \theta$.
  \item Rearrange the terms to group $\csc \theta$ and $\cot \theta$ together.
  \item Use the identity $1 = \csc^2 \theta - \cot^2 \theta$ to substitute the '1' in the denominator.
  \item Factorize the denominator using the difference of squares identity $a^2 - b^2 = (a-b)(a+b)$.
  \item Cancel the common term in the numerator and denominator to arrive at the RHS.
\end{enumerate}
\textbf{Derivation:}
\begin{center}
\fitmath{\begin{aligned}
LHS = \frac{\frac{\cos \theta}{\sin \theta} - \frac{\sin \theta}{\sin \theta} + \frac{1}{\sin \theta}}{\frac{\cos \theta}{\sin \theta} + \frac{\sin \theta}{\sin \theta} - \frac{1}{\sin \theta}} = \frac{\cot \theta - 1 + \csc \theta}{\cot \theta + 1 - \csc \theta} \\
= \frac{\cot \theta + \csc \theta - (\csc^2 \theta - \cot^2 \theta)}{\cot \theta - \csc \theta + 1} \\
= \frac{(\cot \theta + \csc \theta) - [(\csc \theta - \cot \theta)(\csc \theta + \cot \theta)]}{\cot \theta - \csc \theta + 1} \\
= \frac{(\cot \theta + \csc \theta) [1 - (\csc \theta - \cot \theta)]}{\cot \theta - \csc \theta + 1} \\
= \frac{(\cot \theta + \csc \theta) [1 - \csc \theta + \cot \theta]}{\cot \theta - \csc \theta + 1} \\
= \cot \theta + \csc \theta = RHS
\end{aligned}}
\end{center}
\hfill \textit{Introduction to Trigonometry — Trigonometric Identities}
\vspace{0.3cm}

\question \textbf{38.} A tower stands vertically on the ground. From a point on the ground, which is $15$ m away from the foot of the tower, the angle of elevation of the top of the tower is found to be $60^\circ$. The height of the tower is:

\textbf{Answer:}
(C)

\textbf{Solution:}
\begin{enumerate}
  \item Let the height of the tower be $h$.
  \item Given distance from foot = $15$ m.
  \item In $\triangle ABC$, $\tan(60^\circ) = \frac{\text{height}}{\text{base}}$.
  \item $\sqrt{3} = \frac{h}{15}$
  \item $h = 15\sqrt{3}$
\end{enumerate}
\textbf{Derivation:}
\begin{center}
\fitmath{\begin{aligned}
\tan(60^\circ) = \frac{h}{15} \\
\implies \sqrt{3} = \frac{h}{15} \\
\implies h = 15\sqrt{3} \text{ m}
\end{aligned}}
\end{center}
\hfill \textit{Some Applications of Trigonometry — Heights and Distances}
\vspace{0.3cm}

\question \textbf{39.} The angle of elevation of the top of a pole from a point on the ground is $30^\circ$. If the pole is $10$ m high, the distance of the point from the foot of the pole is:

\textbf{Answer:}
(B)

\textbf{Solution:}
\begin{enumerate}
  \item Let the distance be $x$.
  \item Given height = $10$ m, angle = $30^\circ$.
  \item $\tan(30^\circ) = \frac{10}{x}$
  \item $\frac{1}{\sqrt{3}} = \frac{10}{x}$
  \item $x = 10\sqrt{3}$
\end{enumerate}
\textbf{Derivation:}
\begin{center}
\fitmath{\begin{aligned}
\tan(30^\circ) = \frac{10}{x} \\
\implies \frac{1}{\sqrt{3}} = \frac{10}{x} \\
\implies x = 10\sqrt{3}
\end{aligned}}
\end{center}
\hfill \textit{Some Applications of Trigonometry — Heights and Distances}
\vspace{0.3cm}

\question \textbf{40.} If the length of a string of a kite flying at a height of $60$ m from the ground is $120$ m, then the angle of inclination of the string with the ground is:

\textbf{Answer:}
(A)

\textbf{Solution:}
\begin{enumerate}
  \item Let the angle be $\theta$.
  \item Height (opposite) = $60$ m, Hypotenuse = $120$ m.
  \item $\sin(\theta) = \frac{60}{120}$
  \item $\sin(\theta) = \frac{1}{2}$
  \item $\theta = 30^\circ$
\end{enumerate}
\textbf{Derivation:}
\begin{center}
\fitmath{\begin{aligned}
\sin(\theta) = \frac{60}{120} = \frac{1}{2} \\
\implies \theta = 30^\circ
\end{aligned}}
\end{center}
\hfill \textit{Some Applications of Trigonometry — Heights and Distances}
\vspace{0.3cm}

\end{questions}

\end{document}
